\section{Two-domain-wall trial vectors}
\label{sec:uncle_two_wall_trial}

The variational argument of
\cite[equation defining $\phi_N$, low-energy-state lemma]{FernandezGonzalez2015Uncle}
uses the earlier normalized standard form: the two diagonal blocks have
unital transfer operators of spectral radius one, faithful trace-one dual
fixed densities, and no other peripheral eigenvalues. Distinct blocks have
mixed transfer spectral radius strictly below one. These conditions are
needed for the asymptotic estimates and do not follow merely from the full
perturbation span condition in Theorem~\ref{thm:uncle_periodic_ground_space}.

\begin{definition}[The finite two-wall trial]
    \label{def:uncle_two_wall_trial}
    \lean{MPSTensor.uncleTwoWallCoefficient,
        MPSTensor.uncleTwoWallCoefficient_eq_sum,
        MPSTensor.uncleTwoWallTrial,MPSTensor.uncleTwoWallTrialSource}
    \notready
    \uses{def:uncle_boundary_space}
    For five consecutive words $u,v,w,x,y$, define the coefficient
    \[
        \operatorname{tr}\!\left(A(u)\,R(v)\,B(w)\,L(x)\,A(y)\right),
    \]
    where $R(v)$ and $L(x)$ are the sums with one rectangular wall at each
    position of their respective words. Thus the coefficient is precisely
    the independent double sum over both wall positions. On a chain of
    length $6N+1$ with $N\ge1$, choose consecutive segment lengths
    $N,N+1,2N-1,N+1,N$. In coordinates numbered from zero, the first wall
    ranges from $N$ through $2N$, and the second from $4N$ through $5N$.
    Denote the resulting vector by $\phi_N$.
\end{definition}

\begin{lemma}[Interior moving-wall cancellation]
    \label{lem:uncle_two_wall_interior_cancellation}
    \lean{MPSTensor.uncleBoundaryWord_slice_mem,
        MPSTensor.uncleInteractionES_upper_wall_slice_eq_zero,
        MPSTensor.uncleInteractionES_lower_wall_slice_eq_zero,
        MPSTensor.uncleInteractionES_twoWall_slice_first_eq_zero,
        MPSTensor.uncleInteractionES_twoWall_slice_second_eq_zero}
    \notready
    \uses{def:uncle_two_wall_trial,thm:uncle_periodic_ground_space}
    Fix all exterior physical letters, and vary two adjacent physical
    letters entirely inside either moving-wall interval. The resulting
    two-site vector is annihilated by the actual uncle local projector.
    This algebraic cancellation requires neither perturbation injectivity
    nor canonical normalization.
\end{lemma}
\begin{proof}
    Concatenation multiplies the single-wall coefficient words by the
    bilinear product defining $S_n$. A fixed exterior contraction therefore
    composes a linear functional with multiplication by the two fixed
    exterior coefficient words. Trace duality represents that functional
    by a boundary matrix, so its two-site range lies in $S_2$, the actual
    local projector kernel. Reading the upper or lower rectangular block
    gives the first or second wall interval, respectively. For the latter,
    the rectangular trace identity rotates the fixed exterior first wall
    around the trace. Each wall position is included in the moving sum.
\end{proof}

The four interval-end contributions, the asymptotic norm and ground-sector
orthogonality estimates, the proper algebraic exceptional sets, and the
variational spectral conclusion are separate steps. The preceding local
cancellation does not yet assert finite-volume gaplessness.
